\chapter{Cyclosporine Level}

\section*{What Is This Test?}
The cyclosporine level measures the whole blood concentration of cyclosporine (cyclosporin A), a calcineurin inhibitor immunosuppressant used to prevent solid organ transplant rejection (kidney, liver, heart, lung, pancreas) and to treat graft-versus-host disease (GVHD) after hematopoietic stem cell transplantation. Cyclosporine binds to cyclophilin, and the cyclosporine-cyclophilin complex inhibits calcineurin, preventing dephosphorylation and nuclear translocation of nuclear factor of activated T cells (NFAT), thereby blocking interleukin-2 (IL-2) gene transcription and T-cell activation. Cyclosporine is metabolized by CYP3A4/CYP3A5 in the liver and small intestine and is a substrate of P-glycoprotein (ABCB1). It has a narrow therapeutic index: subtherapeutic levels lead to acute rejection and chronic allograft nephropathy; supratherapeutic levels cause acute and chronic nephrotoxicity (afferent and efferent arteriolar vasoconstriction, tubular atrophy, interstitial fibrosis), hypertension, neurotoxicity (tremor, headache, seizures, posterior reversible encephalopathy syndrome/PRES), hyperkalemia, hypomagnesemia, hyperuricemia/gout, gingival hyperplasia, and hirsutism. TDM is essential. \cite{kahan2004therapeutic} The preferred specimen is whole blood (EDTA, lavender-top tube) because cyclosporine is extensively distributed in erythrocytes (50--60\%) and lipoproteins. Trough (C0, pre-dose) or 2-hour post-dose (C2) monitoring strategies are used depending on the transplant center.

\section*{Alternative Names}
Cyclosporine Level; Cyclosporin A; CsA Level; Sandimmune; Neoral; Gengraf; Cyclosporine TDM; Whole Blood Cyclosporine

\section*{Principle}
Cyclosporine is measured by whole blood liquid chromatography-tandem mass spectrometry (LC-MS/MS), the gold standard, or by automated chemiluminescent microparticle immunoassay (CMIA, Abbott Architect), or enzyme multiplied immunoassay (EMIT, Siemens). LC-MS/MS directly quantifies the cyclosporine molecule and is preferred because it does not cross-react with inactive cyclosporine metabolites that accumulate in hepatic dysfunction. Immunoassays overestimate cyclosporine concentrations by 10--30\% due to cross-reactivity with metabolites (AM1, AM9, AM4N) \cite{holt2002international}. Specimen: whole blood collected in EDTA (lavender-top) tube --- serum or plasma are NOT used because cyclosporine is concentrated in erythrocytes and drug concentration is temperature-dependent (redistribution between plasma and RBCs).

\section*{History}
Cyclosporine was discovered in 1970 from the soil fungus Tolypocladium inflatum (originally isolated from a soil sample collected in Norway) by scientists at Sandoz (now Novartis). Its immunosuppressive properties were discovered by Jean-Francois Borel in 1976 \cite{borel1976biological}. The first successful use in human kidney transplantation was reported by Roy Calne in 1978 \cite{calne1978cyclosporin}. FDA approval for transplantation was granted in 1983. Cyclosporine revolutionized solid organ transplantation, dramatically reducing acute rejection rates and improving 1-year graft survival from approximately 50\% (pre-cyclosporine era, azathioprine + prednisone) to 80--90\%.

\section*{Why This Test Is Ordered}
Routine TDM in transplant recipients to maintain cyclosporine levels within target ranges (trough C0 or C2 monitoring, varying by transplant type and time post-transplant). Evaluation of suspected cyclosporine nephrotoxicity versus acute rejection. Dose adjustment after introduction or discontinuation of CYP3A4/P-gp interacting drugs. Monitoring when changing cyclosporine formulations (Sandimmune to Neoral/Gengraf microemulsion). The microemulsion formulation (Neoral, Gengraf) has ~30\% higher bioavailability than Sandimmune and is not interchangeable.

\section*{Primary vs Secondary Test}
TDM is a secondary monitoring test. Cyclosporine dosing is primarily managed by graft function and clinical tolerability; the whole blood trough (C0) or 2-hour post-dose (C2) level is the essential secondary tool for individualizing therapy within transplant-specific targets and for helping to distinguish rejection from calcineurin-inhibitor nephrotoxicity.

\section*{How This Test Is Performed}
A blood sample is collected by standard venipuncture into an EDTA (lavender-top) tube; whole blood is mandatory because cyclosporine is extensively distributed in erythrocytes and lipoproteins, and serum or plasma is not used. Cyclosporine is measured by LC-MS/MS (the gold standard) or by automated immunoassay (CMIA or EMIT), which cross-reacts with inactive metabolites and may overestimate concentrations by 10--30\% \cite{rifai2018tietz}. Results are available within a few hours for urgent requests and 1--2 days for routine samples; levels are reported as a trough (C0, pre-dose) or C2 (2 hours after the dose), depending on the center's protocol.

\section*{What Preparation Is Needed}
The sample must be drawn as a trough (C0, immediately before the next dose) or exactly 2 hours after a witnessed dose for C2 monitoring, at steady state (approximately 3--5 days after initiation or dose change). Dosing should be at consistent 12-hour intervals. Because cyclosporine is a CYP3A4/P-gp substrate, the time of the last dose, the formulation (Sandimmune vs. Neoral/Gengraf microemulsion, which are not interchangeable), and any interacting drugs (azole antifungals, macrolides, calcium channel blockers, rifampin, phenytoin) must be documented for interpretation.

\section*{Contraindications and Precautions}
\begin{itemize}
  \item No absolute contraindications to standard venipuncture
  \item Whole blood EDTA is mandatory; serum, plasma, and heparinized samples are not acceptable
  \item Immunoassay results may overestimate true concentrations in hepatic dysfunction owing to accumulating inactive metabolites; LC-MS/MS is preferred in this setting
  \item Interpretation requires the transplant type, time post-transplant, and monitoring strategy (C0 vs. C2), which are not interchangeable
  \item Nephrotoxicity and hypertension can occur at levels within the target range, and rising creatinine during cyclosporine therapy may reflect drug toxicity or rejection, requiring allograft biopsy to distinguish
  \item Microemulsion formulations (Neoral, Gengraf) have approximately 30\% higher bioavailability than Sandimmune and are not bioequivalent; conversion between formulations requires level monitoring
\end{itemize}

\section*{Risks}
Minimal risk. Slight pain or bruising at the venipuncture site. Rarely: hematoma, infection, or vasovagal syncope.

\section*{What Happens After the Test}
The result is interpreted against the transplant-specific target (C0 or C2, by organ and time post-transplant). Subtherapeutic levels prompt a dose increase with repeat sampling at a new steady state; supratherapeutic levels with rising creatinine prompt dose reduction and rechecking of the level and renal function. Follow-up levels are typically obtained 3--5 days after a dose change and after starting or stopping CYP3A4/P-gp interacting drugs.

\section*{Understanding Results}
Subtherapeutic: Trough < target range. Risk of acute rejection. Supratherapeutic: Trough > target range. Toxicity: acute nephrotoxicity (reversible afferent arteriolar vasoconstriction, rising creatinine), chronic nephrotoxicity (irreversible striped interstitial fibrosis/tubular atrophy), neurotoxicity (tremor, headache, PRES), hypertension, hyperkalemia, hypomagnesemia, gingival hyperplasia, hirsutism \cite{naesens2009calcineurin}.

\section*{Normal Results}
Kidney transplant: C0 (trough) 150--300 ng/mL (first 3 months), 100--200 ng/mL (maintenance >3 months). Liver transplant: C0 200--300 ng/mL (early), 100--150 ng/mL (maintenance). Heart transplant: C0 250--350 ng/mL (first 6 months), 150--250 ng/mL (maintenance). C2 monitoring (2 hours post-dose): kidney transplant C2 800--1,200 ng/mL (early), 600--800 ng/mL (maintenance). Specimen: whole blood EDTA --- essential.

\section*{Follow-up Tests}
Serum creatinine, BUN, eGFR (before and serially), serum electrolytes (potassium, magnesium), blood pressure, uric acid. Concurrent medication review for CYP3A4/P-gp interactions. Allograft biopsy if rejection versus drug toxicity is unclear.

\section*{References}
\bibliographystyle{unsrtnat}
\bibliography{references}

\clearpage
